\section{Síntesis y ampliación}

La lógica de esta clase puede condensarse en una cadena cuya dificultad aumenta en cada eslabón. Primero, la human intent real no se limita a la instrucción literal; segundo, RLHF convierte parte de las preferencias implícitas en una señal optimizable mediante demonstrations, comparisons, un reward model y RL; tercero, InstructGPT y las primeras versiones de ChatGPT demostraron que esta señal mejora de forma notable la utilidad práctica, aunque persisten la hallucination, la prompt sensitivity y los problemas de representatividad; cuarto, cuando las tareas que completa el modelo rebasan la capacidad humana de evaluarlas de forma independiente, la señal del RLHF convencional se distorsiona; quinto, la AI-assisted evaluation busca que la máquina se encargue de leer, buscar, calcular y elaborar la critique, para que el ser humano se concentre en expresar sus preferencias; sexto, las targeted perturbations, los critique experiments y la discriminator--critique gap solo permiten medir paso a paso si esta línea de trabajo realmente amplía la supervisión.

\begin{importantbox}{Los cinco puntos más importantes de toda la clase}
\begin{enumerate}
  \item El objeto de la alignment es la human intent, no una obedience mecánica.
  \item El reward model es un proxy reutilizable de un conjunto limitado de datos de comparación y, al mismo tiempo, una nueva superficie de ataque expuesta a la optimización.
  \item La preference improvement, la capability, la truthfulness y la safety deben medirse por separado.
  \item Lo esencial de la scalable oversight no es que la IA decida los valores en lugar de las personas, sino transformar los objetos difíciles de evaluar en evidencia verificable.
  \item Todas las líneas de supervisión recursiva, self-critique, debate e interpretabilidad terminan enfrentando la misma pregunta: de dónde proviene una ground truth confiable o una trusted evaluation signal.
\end{enumerate}
\end{importantbox}

\subsection{Escalera de madurez de la evidencia}

El lector puede usar una escalera de cuatro niveles para juzgar la solidez de una alignment claim. El nivel de mecanismo explica cómo funciona el método; el nivel experimental actual indica si es eficaz en tareas controladas; el nivel de transferencia comprueba si se aplica a tareas más difíciles, más realistas y más adversariales; solo el nivel de garantía discute si sigue siendo confiable frente a sistemas más capaces que el evaluator. Los resultados de RLHF de esta clase se ubican principalmente en los dos primeros niveles; los resultados de scalable oversight son una proof of concept temprana del segundo nivel, y el papel final de la outer/inner alignment y de la interpretabilidad sigue siendo un tema de investigación abierto.

\begin{warningbox}{No presentar una línea de investigación como un hecho consumado}
La clase planteó una dirección prometedora: usar IA para ayudar a las personas a evaluar IA. No demostró que la curva de evaluación pueda alcanzar siempre a la curva de capacidad, ni que un mismo modelo vaya a revelar con honestidad todo lo que sabe. Unas notas de calidad deben conservar esta incertidumbre, porque la dificultad del problema radica justamente en que no podemos verificar directamente los sistemas potentes del futuro con los experimentos pequeños de hoy.
\end{warningbox}

\subsection{Lecturas adicionales}

Todos los materiales siguientes se publicaron antes de la fecha de la clase y respaldan directamente su línea principal:

\begin{itemize}
  \item \href{https://arxiv.org/abs/1706.03741}{Christiano et al., Deep Reinforcement Learning from Human Preferences}: preferencias por pares de trayectorias y reward learning en línea.
  \item \href{https://arxiv.org/abs/1811.07871}{Leike et al., Scalable Agent Alignment via Reward Modeling}: agenda de investigación de reward modeling y recursive assistance.
  \item \href{https://arxiv.org/abs/2203.02155}{Ouyang et al., Training Language Models to Follow Instructions with Human Feedback}: implementación pública y evaluación de SFT, RM, PPO y PPO-ptx.
  \item \href{https://arxiv.org/abs/2009.01325}{Stiennon et al., Learning to Summarize with Human Feedback}: optimización de preferencias en la generación de lenguaje.
  \item \href{https://arxiv.org/abs/2206.05802}{Saunders et al., Self-critiquing Models for Assisting Human Evaluators}: critique assistance, targeted flaws y discriminator--critique gap.
  \item \href{https://arxiv.org/abs/1805.00899}{Irving et al., AI Safety via Debate}: otra línea que amplifica el juicio humano mediante el diálogo y la argumentación adversarial.
\end{itemize}

En última instancia, la alignment no consiste en conectar un «goodness score» a la salida del modelo, sino en diseñar una cadena de retroalimentación que no pierda confiabilidad a medida que crece la capacidad. Al final, el ponente no prometió que algún componente resolvería todos los problemas; la postura de investigación más confiable que propuso es: hacer medibles las tareas de evaluación, hacer verificable la evidencia, hacer explícitos los límites de falla, y seguir preguntando qué sabe el modelo, qué dice y qué hace en realidad.

\subsection{Autoevaluación posterior a la clase}

Después de leer esta clase, no basta con saber repetir el flujo de RLHF; también hay que poder determinar qué tramo de la cadena de retroalimentación mejora realmente una nueva propuesta de alignment. Las cuatro preguntas siguientes sirven como lista mínima de verificación al leer artículos, diseñar experimentos o revisar las afirmaciones de un producto.

\begin{knowledgebox}{Cuatro preguntas para examinar una nueva alignment claim}
\begin{enumerate}
  \item ¿La señal de entrenamiento proviene de demonstrations, comparisons, verifiable outcomes o del juicio de otro modelo no validado?
  \item ¿El experimento mide la preference, la tasa real de aciertos, la calibration, la robustness, o solo un proxy fácil de optimizar?
  \item Cuando la dificultad de la tarea rebasa la capacidad humana de evaluación independiente, ¿qué evidence verificable aporta el sistema, más allá de una explicación más fluida?
  \item Cuando ocurre una falla, ¿es posible distinguir entre que el modelo no sabe, que sabe pero no puede expresarlo, que sabe pero no quiere decirlo, y que el evaluator no lo detectó?
\end{enumerate}
\end{knowledgebox}

\end{document}
